\section{Общая информация о микромеханических датчиках}
\label{sec:general_info_nt}

Микромеханические датчики (микроэлектромеханические системы, МЭМС) — это миниатюрные устройства, которые сочетают механические элементы и электронные схемы. Размеры компонентов таких датчиков варьируются от 1 до 100 микрометров, а самих устройств — от нескольких миллиметров до нескольких сантиметров \cite{adams_layton_2010}.


МЭМС-устройства обычно изготавливают на кремниевой подложке с помощью технологии микрообработки, аналогично технологии изготовления однокристальных интегральных микросхем \cite{adams_layton_2010}.

В основе их работы лежат различные физические законы, благодаря которым можно преобразовать различные механические воздействия в электрический сигнал и наоборот.


МЭМС (микроэлектромеханические системы) датчики широко применяются в различных отраслях благодаря их миниатюрным размерам, низкому энергопотреблению и высокой точности. В качестве основных областей применения можно выделить:

\begin{itemize}
	\item бытовая электроника;
	\item промышленность;
	\item медицина и биотехнологии;
	\item энергетика;
\end{itemize}

Для МЕМС датчиков бытового назначения основным требованием является высокая энергоэфективность, малые габаритные размеры, а также низкая стоимость производства. В то же время, датчики, которые используются в более высокоточных системах должны обладать низкими шумовыми характеристики и высокой чувствительностью.
